% =================================================================
%  style/highlight.tex
%  Палитра хайлайтов и команды выделения.
%
%  В тексте:      \hlY{текст}   — жёлтый маркер
%                 \hlB \hlG \hlR \hlO \hlV — синий, зелёный, красный,
%                 оранжевый, фиолетовый
%                 \hlc{hlBlue}{текст} — произвольный цвет палитры
%  В формуле:     $a + \hlR{b} = c$  — команды сами переключают режим
%  В tikz:        \draw[hlRed, hlcurve] ...;   \fill[hlGreen, hlarea] ...;
%
%  Выделение рассчитано на короткие фрагменты внутри строки: заливка
%  не переносится на следующую строку (это ограничение \colorbox).
%  Для длинного куска разбейте его на несколько \hl... или выделите
%  абзац боксом.
% =================================================================

% -----------------------------------------------------------------
% Палитра
% -----------------------------------------------------------------
\definecolor{hlYellow}{RGB}{250, 220, 130}
\definecolor{hlBlue}{RGB}{155, 190, 220}
\definecolor{hlGreen}{RGB}{160, 195, 170}
\definecolor{hlRed}{RGB}{220, 155, 160}
\definecolor{hlOrange}{RGB}{225, 175, 125}
\definecolor{hlViolet}{RGB}{185, 170, 205}

% -----------------------------------------------------------------
% Базовая обёртка.
%   \hlc{<цвет>}{<содержимое>}
% \fboxsep уменьшен, чтобы заливка облегала текст, а не висела рамкой.
% В математическом режиме содержимое возвращается в математику через
% \text{...}{$...$}, иначе \colorbox сбросил бы его в текстовый режим.
% -----------------------------------------------------------------
\newlength{\hlsep}
\setlength{\hlsep}{1.5pt}

\newcommand{\hlc}[2]{%
	\begingroup
	\setlength{\fboxsep}{\hlsep}%
	\ifmmode
		\text{\colorbox{#1}{$#2$}}%
	\else
		\colorbox{#1}{#2}%
	\fi
	\endgroup
}

% -----------------------------------------------------------------
% Короткие команды по цветам
% -----------------------------------------------------------------
\newcommand{\hlY}[1]{\hlc{hlYellow}{#1}}
\newcommand{\hlB}[1]{\hlc{hlBlue}{#1}}
\newcommand{\hlG}[1]{\hlc{hlGreen}{#1}}
\newcommand{\hlR}[1]{\hlc{hlRed}{#1}}
\newcommand{\hlO}[1]{\hlc{hlOrange}{#1}}
\newcommand{\hlV}[1]{\hlc{hlViolet}{#1}}

% -----------------------------------------------------------------
% Стили для иллюстраций: цвет берётся из палитры отдельным ключом,
%   \draw[hlRed, hlcurve] plot ...;   — выделенная кривая
%   \fill[hlGreen, hlarea] ...;       — подсвеченная область
%   \node[hlBlue, hlmark] at (x,y) {}; — отмеченная точка
% -----------------------------------------------------------------
\tikzset{
	hlcurve/.style = {line width=1.2pt, line cap=round},
	hlarea/.style  = {opacity=0.35},
	hlmark/.style  = {circle, fill, inner sep=1.6pt},
}
